\section{Getting started}

\subsection{Building the CLI}

Build the customer CLI with a release build; the binary is
\file{target/release/keyquorum}.

\begin{lstlisting}
cargo build --release
\end{lstlisting}

Adding \flag{--features provider} compiles mailbox-host capabilities into
the same crate. That flag is a build capability, not authorization: official
clients trust a relay only when it presents a KeyQuorum-signed provider
certificate and proves possession of the matching private key. This manual
covers the customer-facing surface --- the URL and API key a provider gives
you, and the \code{loadkey} / \code{relay push} / \code{relay pull} commands
built on top of them. The rest of this manual explains what each area of
the CLI is for and how the pieces fit together; the
\hyperref[sec:command-reference]{Command reference} at the end is the
complete index of every \code{keyquorum} and \code{keyquorum-device}
command and subcommand. \code{keyquorum <command> --help} still has the
full, current flag list for any one command.

\subsection{Quick start}

Most tasks are one command once you have an identity. \code{setup} creates it
and \code{use} remembers it, so later commands can leave \flag{--as},
\flag{--slot} and \flag{--url} out; a flag you type always wins.

\begin{lstlisting}
keyquorum setup --device ./usb --label alice --url https://relay.example.com
keyquorum doctor --to bob
keyquorum send report.pdf --to bob
keyquorum inbox
keyquorum inbox open
\end{lstlisting}

\code{setup} creates the device container if it is missing, provisions your
slot, registers its encryption and signing keys, binds the slot, loads the relay
key and runs \code{use}. It can be run again; a step that is already done is
skipped. \code{doctor} reads only and prints, for each missing thing, the
command that fixes it. \code{use --show} prints the stored defaults, which are
pointers only (label, slot, container path, relay URL, cache setting): no
passphrase, key or bearer is stored there.

A left-out flag resolves, in order, from the flag you typed, a recent
parameter, the stored default, and then the command's own behaviour or error.
The environment variable \code{KEYQUORUM\_DB} chooses the database when
\flag{--db} is not given.

The older, longer procedures --- \code{deliver send|open|ack}, \code{file
share|receive|ack} and \code{relay pull} --- are still possible and behave as
they always did, but they are not recommended: each prints a note naming its
replacement. Section~\ref{sec:legacy-commands} lists them.

Three short-lived caches are fresh for 15 minutes: parameters you just used (for
the same command only, and never a recipient or any target of an outward or
destructive command), a relay identity check that just passed (never a failure,
and \code{loadkey} always runs the full check), and facts \code{doctor} already
looked up. Every reuse is announced on standard error. Turn them off with
\flag{--no-cache}, \code{KEYQUORUM\_NO\_CACHE=1} or \code{use --cache off};
\code{cache clear} empties them.

\subsection{The local database}

Most commands read and write a local SQLite database that holds your split
trees, registered keys, and (for a relay-connected identity) a sealed copy
of your mailbox bearer. Pass \flag{--db} to point at a specific file; when
omitted, commands use \code{KEYQUORUM\_DB} if it is set, and otherwise a
default database in the current directory. Running
the same commands against different \flag{--db} files is how one machine
can hold several independent identities (for example \file{alice.sqlite}
and \file{bob.sqlite} in the walkthroughs throughout this manual).

\begin{securityBox}
Never commit a SQLite database, a private key file, an API bearer, a
provider certificate, a revocation list, a device descriptor
(\file{device.kq}, \file{device.skey}), or a slot token (\file{*.kqst}) to
version control. See \file{.gitignore} for the patterns this repository
already excludes.
\end{securityBox}

\subsection{Two kinds of identity}

KeyQuorum draws a line between two things that both use the word ``key'':

\begin{itemize}
  \item A \textbf{hardware key} is a registered Ed25519/X25519 keypair,
  generated with \code{keyquorum generate} or brought in with
  \code{keyquorum register}. It is what a person, device, or slot presents
  to satisfy a quorum, sign a message, or open a sealed envelope.
  \item A \textbf{split key} (or ``key'' in \code{split} / \code{tree}
  output) is the secret being protected --- a standalone escrow, or the
  data key encrypting a specific file. It lives as a tree of leaves, each
  leaf sealed to one hardware key.
\end{itemize}

The rest of this manual moves through the CLI roughly in the order you
would use it for a new deployment: registering hardware keys and devices,
building a split tree, protecting files and credentials against it, then
the supporting machinery --- private sign bridges, authenticated updates,
device transfer, direct file delivery, and the mailbox relay that connects
machines that are not open together.
